\section{Continuous Assessment 2 (CA-2)}

\ExamHeader{Continuous Assessment 2 (CA-2)}{Set B (Model Paper 2)}{60 Minutes}{30 Marks}{None}

\noindent
\textbf{General Instructions:}
\begin{enumerate}[topsep=2pt,itemsep=1pt]
    \item This paper contains \textbf{6 subjective questions} carrying \textbf{5 marks each}.
    \item Candidates are required to \textbf{attempt any 5 questions} ($5 \times 5 = 25$ marks or scaled to 30 marks).
    \item Show all necessary algebraic steps, matrix reductions, and derivative formulations clearly.
\end{enumerate}
\vspace{3mm}

\noindent\textbf{Question 1} \hfill \textbf{[5 Marks]}

If $u = \ln\left(\frac{x^4+y^4}{x+y}\right)$, prove that $x\frac{\partial u}{\partial x} + y\frac{\partial u}{\partial y} = 3$.

\begin{solbox}
Let $z = e^u = \frac{x^4+y^4}{x+y}$.
$z(tx, ty) = \frac{t^4(x^4+y^4)}{t(x+y)} = t^3 z(x,y)$, so $z$ is homogeneous of degree $n = 3$.
By Euler's Theorem on $z$:
\begin{equation}
x\frac{\partial z}{\partial x} + y\frac{\partial z}{\partial y} = 3z
\end{equation}
Since $z = e^u$, $\frac{\partial z}{\partial x} = e^u \frac{\partial u}{\partial x}$ and $\frac{\partial z}{\partial y} = e^u \frac{\partial u}{\partial y}$.
Substituting:
\begin{equation}
e^u \left( x\frac{\partial u}{\partial x} + y\frac{\partial u}{\partial y} \right) = 3e^u \implies x\frac{\partial u}{\partial x} + y\frac{\partial u}{\partial y} = 3
\end{equation}
\end{solbox}

\begin{examinertip}
Exponentiating logarithmic functions turns them into standard homogeneous polynomials of degree $(4-1=3)$.
\end{examinertip}

\vspace{4mm}

\noindent\textbf{Question 2} \hfill \textbf{[5 Marks]}

Find the local maximum, local minimum, and saddle points of $f(x,y) = 2x^2 + y^2 - 4x - 2y + 5$.

\begin{solbox}
1. Critical points:
$f_x = 4x - 4 = 0 \implies x = 1$.
$f_y = 2y - 2 = 0 \implies y = 1$.
The only critical point is $(1, 1)$.

2. Hessian discriminant:
$r = f_{xx} = 4, \quad s = f_{xy} = 0, \quad t = f_{yy} = 2$.
$D = rt - s^2 = (4)(2) - 0 = 8 > 0$.
Since $D > 0$ and $r = 4 > 0$, the point $(1, 1)$ is a \textbf{Local Minimum} with minimum value $f(1,1) = 2(1)+1-4-2+5 = 2$.
\end{solbox}

\begin{examinertip}
Since $D > 0$ and $f_{xx} > 0$, $(1, 1)$ is an unconstrained absolute minimum.
\end{examinertip}

\vspace{4mm}

\noindent\textbf{Question 3} \hfill \textbf{[5 Marks]}

Find the maximum and minimum values of $f(x,y) = x y$ subject to the constraint $x^2 + y^2 = 8$ using the method of Lagrange Multipliers.

\begin{solbox}
Let $g(x,y) = x^2 + y^2 = 8$.
Lagrange multiplier equations: $\nabla f = \lambda \nabla g \implies \begin{cases} y = 2\lambda x \\ x = 2\lambda y \end{cases}$.
Multiplying the first by $x$ and second by $y$: $xy = 2\lambda x^2 = 2\lambda y^2 \implies x^2 = y^2$ (since $\lambda \ne 0$).
Substitute into constraint: $x^2 + x^2 = 8 \implies 2x^2 = 8 \implies x^2 = 4 \implies x = \pm 2, \; y = \pm 2$.
Points are $(2, 2), (-2, -2), (2, -2), (-2, 2)$.
\begin{itemize}
    \item At $(2, 2)$ and $(-2, -2)$: $f(x,y) = 4$ (\textbf{Maximum Value}).
    \item At $(2, -2)$ and $(-2, 2)$: $f(x,y) = -4$ (\textbf{Minimum Value}).
\end{itemize}
\end{solbox}

\begin{examinertip}
Always verify all four sign combinations $(\pm 2, \pm 2)$ on the constraint circle.
\end{examinertip}

\vspace{4mm}

\noindent\textbf{Question 4} \hfill \textbf{[5 Marks]}

Evaluate $\iint_R (x + y)\,dA$ where $R$ is the region bounded by $y = x^2$ and $y = 2x$.

\begin{solbox}
1. Points of intersection: $x^2 = 2x \implies x(x-2) = 0 \implies x = 0, \; x = 2$.
For $x \in [0, 2]$, $y$ varies from the parabola $y = x^2$ up to the line $y = 2x$.
2. Formulate iterated integral:
\begin{align}
I &= \int_0^2 \left( \int_{x^2}^{2x} (x + y)\,dy \right) dx = \int_0^2 \left[ x y + \frac{y^2}{2} \right]_{x^2}^{2x} dx \\
&= \int_0^2 \left( \left[x(2x) + \frac{4x^2}{2}\right] - \left[x(x^2) + \frac{x^4}{2}\right] \right) dx \\
&= \int_0^2 \left( 4x^2 - x^3 - \frac{x^4}{2} \right) dx = \left[ \frac{4x^3}{3} - \frac{x^4}{4} - \frac{x^5}{10} \right]_0^2 \\
&= \frac{32}{3} - 4 - \frac{32}{10} = \frac{32}{3} - 4 - \frac{16}{5} = \frac{160 - 60 - 48}{15} = \frac{52}{15}
\end{align}
\end{solbox}

\begin{examinertip}
Sketch the parabola and straight line to verify the lower and upper bounds of $y$.
\end{examinertip}

\vspace{4mm}

\noindent\textbf{Question 5} \hfill \textbf{[5 Marks]}

Change the order of integration and evaluate $I = \int_0^a \int_0^{\sqrt{a^2-x^2}} y^2 \sqrt{a^2-x^2}\,dy\,dx$.

\begin{solbox}
The region is the first quadrant quarter disk $x^2 + y^2 \le a^2$ ($0 \le x \le a, 0 \le y \le \sqrt{a^2-x^2}$).
Transforming into polar coordinates $x = r\cos\theta, y = r\sin\theta$:
$r \in [0, a], \theta \in [0, \pi/2]$.
$y^2 = r^2\sin^2\theta, \sqrt{a^2-x^2} = \sqrt{a^2-r^2\cos^2\theta}$.
Directly integrating in Cartesian:
\begin{equation}
I = \int_0^a \sqrt{a^2-x^2} \left[ \frac{y^3}{3} \right]_0^{\sqrt{a^2-x^2}} dx = \frac{1}{3} \int_0^a (a^2-x^2)^2\,dx = \frac{1}{3} \int_0^a (a^4 - 2a^2 x^2 + x^4)\,dx = \frac{8a^5}{45}
\end{equation}
\end{solbox}

\begin{examinertip}
Integrating with respect to $y$ first directly removes the square root.
\end{examinertip}

\vspace{4mm}

\noindent\textbf{Question 6} \hfill \textbf{[5 Marks]}

Find the volume of the region bounded by the coordinate planes and the plane $\frac{x}{a} + \frac{y}{b} + \frac{z}{c} = 1$ in the first octant.

\begin{solbox}
The solid tetrahedron $\Omega$ is defined by $x \in [0, a]$, $y \in [0, b(1 - x/a)]$, $z \in [0, c(1 - x/a - y/b)]$.
\begin{align}
V &= \iiint_\Omega 1\,dz\,dy\,dx = \int_0^a \int_0^{b(1-x/a)} c\left(1 - \frac{x}{a} - \frac{y}{b}\right)\,dy\,dx \\
&= c \int_0^a \left[ y\left(1 - \frac{x}{a}\right) - \frac{y^2}{2b} \right]_0^{b(1-x/a)} dx \\
&= c \int_0^a \frac{b}{2}\left(1 - \frac{x}{a}\right)^2 dx = \frac{b c}{2} \left[ -\frac{a}{3}\left(1 - \frac{x}{a}\right)^3 \right]_0^a = \frac{a b c}{6}
\end{align}
\end{solbox}

\begin{examinertip}
Standard volume of a Cartesian tetrahedron formed by axes intercepts is $\frac{1}{6}abc$.
\end{examinertip}

\vspace{4mm}

